\chapter{Einleitung}
\label{chap:Einleitung}

\section{Handbuch}
\label{sec:Handbuch}
%% Das Label sollte etwas über den Typ aussagen. Section-Labels sollten mit sec, Chapters mit chap, Bilder mit fig und Tabellen mit tab beginnen. So kann fref den Typ erkennen und entsprechend referenzieren.
Der Reiseführer \gls{pag} hat eine sehr schöne Theorie zur Komplexität des Universums:

\begin{quote}
Es gibt eine Theorie, die besagt, wenn jemals irgendwer genau herausfindet, wozu das Universum da ist und warum es da ist, dann verschwindet es auf der Stelle und wird durch noch etwas Bizarreres und Unbegreiflicheres ersetzt. - Es gibt eine andere Theorie, nach der das schon passiert ist. \hfill \cite{adams1995hitchhiker}
\end{quote}

Sie ist natürlich nicht so schön wie das Logo unserer Uni in \Fref{fig:uni-logo} oder das Logo der \gls{dcsec} in \Fref{fig:dcsec-logo}.

\begin{figure}[ht]
\centering
\includegraphics[width=\textwidth]{logo_luh}
\caption[Uni-Logo]{Das Uni-Logo}
\label{fig:uni-logo}
\end{figure}

\begin{figure}[ht]
\centering
\includegraphics{logo_dcsec}
\caption[DCSec-Logo]{Distributed Computing Security Group Logo}
\label{fig:dcsec-logo}
\end{figure}

\begin{table}
\begin{center}
\begin{tabular}{lcr}
  Spalte 1 & Spalte 2 & Spalte 3 \\
  10 & 20 & 30 \\
\end{tabular}
\end{center}
\label{tab:test}
\caption{Testtabelle}
\end{table}

Was man vielleicht noch beachten sollte ist, dass sind Abkürzungen wie \gls{dcsec} oder \gls{pag} bereits einmal benutzt worden. Werden sie zukünftig nicht mehr neu ausgeschrieben, im Gegensatz zu unbenutzten wie \gls{xml}